\section{Related Work}
Prior work on diffusion model distillation primarily focuses on predicting shortcuts towards less noisy states, with training objectives ranging from direct regression to distribution matching.

Early work~\citep{luhman2021knowledgedistillationiterativegenerative} directly regresses the teacher's ODE integral in a single step, but suffers from degraded quality, since regressing $\vx_0$ with an $\ell_2$ loss tends to produce blurry results. Progressive distillation methods~\citep{salimans2022progressive,liu2022flow,instaflow,frans2025one} make further improvements via a multi-stage process that progressively increases the student's step size and reduces its NFE by regressing the previous stage's multi-step outputs with fewer steps, yet this introduces error accumulation.

Consistency-based models~\citep{consistencymodels,boot,kim2024consistency,song2024improved,meanflow,boffi2025flow} implicitly impose a velocity-based regression loss, which improves quality compared to $\vx$-based regression. However, the flow velocity of a shortcut-predicting student must be constructed implicitly using either inaccurate finite differences or expensive Jacobian--vector products (JVPs). Moreover, their quality is still limited due to accumulation of velocity errors into the integrated shortcut. Therefore, in practice, consistency distillation is often augmented with additional objectives to improve quality~\citep{hypersd,zheng2025largescalediffusiondistillation}, further complicating training.

Conversely, distribution matching approaches~\citep{yin2024onestep,yin2024improved,sauer2025adversarial,sid,sim,salimans2024multistep,imm} adopt score matching and adversarial training to align the student's output distribution with the teacher's. The VSD objective achieves superior quality but risks diversity loss due to mode collapse; GAN and SiD objectives balance quality and diversity but can cause style drift. Their common reliance on auxiliary networks introduces additional tuning complexity and may lead to stability issues at scale~\citep{senseflow}.